\section*{Chapitre \ref{ch:rc:funding-margin-and-capital-adjustments} --- Ajustements de financement, de marge et de capital}

\begin{solution}{exo:rc:funding-margin-and-capital-adjustments:1}
$0.006\times5\,000\,000\times4 = \text{USD}~120\,000$.
\end{solution}

\begin{solution}{exo:rc:funding-margin-and-capital-adjustments:2}
Durée prudentielle $(1-e^{-0.25})/0.05 = 4.42$~; majoration $0.5~\%\times50\text{ millions}\times4.42 =
\text{USD}~1.106$ millions~; sans coût de remplacement, EAD $= 1.4\times1.106 = \text{USD}~1.55$ millions.
\end{solution}

\begin{solution}{exo:rc:funding-margin-and-capital-adjustments:3}
L'\omterm{def:rc:counterparty-exposure:profiles}{exposition négative attendue} (ENE) de la banque dépasse son exposition positive~: sur la courbe à pente ascendante du chapitre,
le receveur de taux fixe devrait verser davantage qu'il ne reçoit sur la durée de vie du swap. Le
collatéral qu'il reçoit au titre de la couverture dans ces états vaut davantage (FBA) que le
collatéral qu'il verse dans les autres états coûte (FCA).
\end{solution}

\begin{solution}{exo:rc:funding-margin-and-capital-adjustments:4}
Les deux valorisent le spread de crédit propre de la banque appliqué à ce qu'elle doit~: la DVA par la probabilité de ne pas payer en cas de défaut,
la FBA par le financement à un taux qui intègre ce spread. Les compter tous deux revient à compter deux fois le spread~;
les banques retirent soit la DVA de la valorisation, soit calculent la FBA sur le spread de financement
net du spread de crédit, ou retiennent une seule FVA symétrique, la DVA étant retirée.
\end{solution}

\begin{solution}{exo:rc:funding-margin-and-capital-adjustments:5}
La CVA facture la perte attendue~; la KVA facture un rendement de 10~\% sur le capital dimensionné pour des pertes situées très loin
dans la queue de distribution (exigence de fonds propres au titre du risque de contrepartie et exigence au titre du risque de CVA), détenu sur toute la durée de vie du swap. Pour une contrepartie notée BBB,
le capital représente plusieurs fois la perte attendue.
\end{solution}

\begin{solution}{exo:rc:funding-margin-and-capital-adjustments:6}
La marge initiale est exigée par la réglementation (règles de marge pour les dérivés non compensés) ou par la chambre de compensation,
doit être ségréguée et ne peut être réutilisée~; elle rapporte moins que ce que la banque paie pour la financer. Elle peut être
réduite en compensant la marge entre transactions (modèles de portefeuille), par la compensation centrale ou par la compression.
\end{solution}

\begin{solution}{exo:rc:funding-margin-and-capital-adjustments:7}
Sans CSA, 63\,700~USD à 40 points de base et 191\,101~USD à 120 points de base~; avec le CSA, 1\,551~USD et 4\,654~USD.
La FCA est linéaire par rapport au spread lorsque celui-ci est plat.
\end{solution}

\begin{solution}{exo:rc:funding-margin-and-capital-adjustments:8}
La couverture du desk crée le besoin de financement, si bien que le coût revient à la transaction~; si le desk
fixe ses prix au taux overnight, la trésorerie subventionne les transactions non collatéralisées et la banque
en traite un volume excessif. Un prix de transfert de fonds qui impute à chaque transaction son usage de financement, net de ce
qu'elle libère, place le coût là où la décision est prise.
\end{solution}

\begin{solution}{pb:rc:funding-margin-and-capital-adjustments:1}
\textbf{1.} FCA de 127\,400~USD, FBA de 185\,797~USD~; FVA nette, un bénéfice de 58\,397~USD.
\textbf{2.} 1.5 points de base par an.
\textbf{3.} 3\,102~USD, 98~\% éliminés.
\textbf{4.} Le client ne verse rien~; la couverture est collatéralisée, si bien que les besoins de trésorerie de la banque suivent
la marge de la couverture, qui reflète la valeur de la transaction client.
\textbf{5.} FCA seule~: 127\,400~USD.
\textbf{6.} CVA de 272\,251~USD et DVA de 210\,351 sans le CSA~; 33\,814 et 17\,745 avec celui-ci.
\textbf{7.} 5.51 millions de dollars sans, 2.35 millions avec.
\textbf{8.} 558\,718~USD sans, 156\,573~USD avec.
\textbf{9.} Le capital impose un plancher sous la marge~: la majoration, ajustée par la période de marge en risque, et le coût de remplacement égal au seuil augmenté du MTA
servent de base au calcul de l'EAD, sur laquelle l'exigence au titre de la CVA est calculée. Le CSA réduit
la KVA de 72~\%, contre 88~\% pour la CVA.
\textbf{10.} MVA de 175\,217~USD, le financement de 4.22 millions de dollars de marge initiale qui décroît
au fil de la durée de vie.
\textbf{11.} 562\,220~USD sans, 347\,595 avec le CSA et la marge initiale.
\textbf{12.} 6.8 et 4.2 points de base par an.
\textbf{13.} Jusqu'à concurrence de la différence, 2.6 points de base par an (4.7 si aucune marge initiale n'était exigée
avec un total de 172\,378~USD, soit 2.1 points de base).
\textbf{14.} Marge de variation quotidienne en espèces ou en obligations éligibles, et éventuellement marge initiale~; le coût
est le financement propre de ce collatéral et les opérations nécessaires pour honorer les appels de marge.
\textbf{15.} Pour cette paire, oui~: un CSA sans marge initiale supprime l'essentiel des coûts de crédit, de financement et
de capital, sans MVA. Les règles qui imposent la marge initiale s'appliquent aux grandes contreparties financières
et, en général, pas aux entreprises non financières.
\textbf{16.} La vision actionnaire pour la cotation du desk (qui paie le financement), la vision en valeur d'entreprise pour la fair value que présentent les comptes,
d'où les conventions du chapitre~20.
\textbf{17.} Parce que le capital est constitué de fonds propres coûteux, qui doivent rapporter le \omterm{def:rc:funding-margin-and-capital-adjustments:kva}{taux de rendement minimal exigé} (hurdle rate)~;
une transaction qui ne rémunère pas son capital détruit de la valeur pour les actionnaires, même rentable avant prise en compte du capital.
\textbf{18.} Par le spread de financement (FCA, FBA, MVA), la survie de la banque dans les pondérations de premier défaut,
et la DVA.
\textbf{19.} Résultat nommé~: \emph{la charge de financement}~: FCA de 127\,400~USD sur le swap décennal
non collatéralisé à 80 points de base, dont le CSA supprime 98~\%.
\textbf{20.} Des charges de crédit, de financement et de capital moindres, contrebalancées
par des coûts opérationnels, de liquidité et de marge.
\end{solution}

\begin{solution}{iq:rc:funding-margin-and-capital-adjustments:1}
La banque couvre le swap non collatéralisé du client par un swap collatéralisé. Lorsque le client
doit à la banque, celle-ci verse une marge sur la couverture et l'emprunte à son taux de financement~; lorsque la
banque doit au client, elle reçoit une marge sur la couverture. La FVA est le coût attendu du premier cas,
diminué du bénéfice attendu du second.
\iqlookfor{la marge de la couverture comme source de financement.}
\end{solution}

\begin{solution}{iq:rc:funding-margin-and-capital-adjustments:2}
Vision en valeur d'entreprise~: les coûts de financement reflètent le risque de défaut de la banque, qui est un transfert aux créanciers, si bien que la
valeur de l'entreprise ne change pas et les prix ne devraient pas inclure la FVA. Vision actionnaire~: ce
transfert est un coût pour les actionnaires, si bien qu'un desk qui maximise la valeur pour l'actionnaire doit le facturer. Une
réponse défendable sépare la fair value de la tarification du desk.
\iqlookfor{les deux points de vue et une position claire.}
\end{solution}

\begin{solution}{iq:rc:funding-margin-and-capital-adjustments:3}
Coût de remplacement $\max(V-C,0)$~; majoration $0.5~\%\times$ notionnel $\times$ durée prudentielle
$\times$ facteur de maturité $\times$ delta, agrégée sur les tranches de maturité de sa devise~;
multiplicateur égal à 1 sauf en cas de surcollatéralisation ou hors de la monnaie~; EAD $= 1.4(\mathrm{RC}+\mathrm{PFE})$.
\iqlookfor{RC, composantes de la majoration et alpha.}
\end{solution}

\begin{solution}{iq:rc:funding-margin-and-capital-adjustments:4}
Le coût de financement de la marge initiale ségréguée sur la durée de vie de la transaction, au spread de financement
au-dessus de ce que rapporte la marge. On le réduit par la marge en portefeuille, en traitant avec des contreparties dont le risque
se compense avec celui de la banque, par la compensation centrale, la compression, et la remise de titres moins coûteux à financer que des espèces.
\iqlookfor{la définition et les leviers.}
\end{solution}

\begin{solution}{iq:rc:funding-margin-and-capital-adjustments:5}
Cela dépend d'une réglementation susceptible d'évoluer, du \omterm{def:rc:funding-margin-and-capital-adjustments:kva}{taux de rendement minimal exigé} (choix de gestion) et de la
projection du capital sur plusieurs décennies~; il peut y avoir double comptage avec d'autres charges. Pour une nouvelle
transaction~: projeter le risque de contrepartie, la CVA et le capital de risque de marché, avec et sans elle, selon les scénarios, puis
prendre la différence et appliquer le \omterm{def:rc:funding-margin-and-capital-adjustments:kva}{taux de rendement minimal exigé} à sa trajectoire actualisée.
\iqlookfor{les doutes et la méthode incrémentale.}
\end{solution}

\begin{solution}{iq:rc:funding-margin-and-capital-adjustments:6}
Des scénarios communs et un moteur d'exposition unique pour tous les \omterm{def:rc:counterparty-exposure:netting}{ensembles de compensation}~; profils stockés une fois et réutilisés par
la CVA, la DVA, la FCA, la FBA et la KVA~; capital réglementaire calculé à partir de la même représentation des transactions~;
différentiation algorithmique pour les sensibilités~; parallélisation sur les trajectoires et les \omterm{def:rc:counterparty-exposure:netting}{ensembles de compensation}~; calculs
incrémentaux pour les nouvelles transactions.
\iqlookfor{scénarios communs, réutilisation et AAD.}
\end{solution}
